\subsubsection{Structural dependence of the electronic coordinate}

The electronic equation expresses a hierarchical composition. Its first
factor, \(\sqrt3/4\), belongs to the central matrix structure and fixes
a geometric normalization. The exponent \(\varphi/\pi^2\) uses the
already generated regional coordinates in the indicated orientation of the
areal--volumetric reader. The term \(22\alpha^3\) incorporates the dodecaphasic
coupling into the cubic correction. The factor \(1+15\Delta_4\) preserves
the global torsional normalization. Finally, the quotient of the action
sections acts with the exponent determined by the return readers.
Each operation thus has a different antecedent within the same
construction.

The articulation becomes more precise when its dependencies are considered.
The central norm can be calculated before the regional coordinates
are evaluated. The exponential and cubic correction require those
coordinates and the closure of \(\alpha\). Multiplicative memory additionally
requires the action sections and the \(108\)-position trajectory.
Knowing the basal expression therefore does not amount to having evaluated
the complete coordinate. The succession of operations does not add an
external mass to the point \((5,5)\); it develops the readers assigned to its
persistent section.

The distinction between scale and structure also provides an
interpretation of the physical comparison. Mass is the dimensional
realization of the complete coordinate, while the centre, involution,
and vacancy describe the relations that produce it. A common change
of unit modifies the dimensional coordinates, but leaves the dimensionless
ratios and the proved transport identities unchanged.
This is the sense in which the electron preserves its structural character
when expressed relative to a chronological scale.

